\section*{第 \ref{ch:b3:banach} 章 --- 巴拿赫空间与
基本定理}

\begin{solution}{exo:b3:banach:1}
(a) $\norm{Sx}_2 = \norm x_2$：$S$ 为等距算子，$\vertiii S =
1$。对后移算子：$\norm{S^*x}_2^2 = \sum_{n \geq
2}\abs{x_n}^2 \leq \norm x_2^2$，等号在 $x = e_2$ 时取到：
$\vertiii{S^*} = 1$。

(b) $\norm{M_ax}_2^2 = \sum\abs{a_n}^2\abs{x_n}^2 \leq \norm
a_\infty^2\norm x_2^2$；取 $x = e_n$ 检验，得 $\vertiii{M_a}
\geq \abs{a_n}$ 对每个 $n$ 成立，故 $\vertiii{M_a} = \norm a_\infty$。

(c) $\abs{\Lambda(f)} \leq \int_0^1\abs f \leq \norm f_\infty$：
$\norm\Lambda \leq 1$。对 $\varepsilon > 0$，令 $f_\varepsilon$
取 $1$ 于 $[0, \frac12 - \varepsilon]$ 上，$-1$ 于 $[\frac12 +
\varepsilon, 1]$ 上，两者之间为仿射：$\norm{f_\varepsilon}_\infty =
1$，且 $\Lambda(f_\varepsilon) \geq 1 - 2\varepsilon$：
$\norm\Lambda = 1$。该上确界不能取到：若 $\Lambda(f) = 1$，且 $\norm
f_\infty \leq 1$，则必有 $\int_0^{1/2}f = \frac12$，且
$\int_{1/2}^1 f = -\frac12$，即（由连续性及 $\abs f \leq 1$）
$f \equiv 1$ 于 $[0, \frac12]$，$f \equiv -1$ 于 $[\frac12,
1]$ 上，在 $\frac12$ 处矛盾。
\end{solution}

\begin{solution}{exo:b3:banach:2}
写 $S = T\bigl(I - T^{-1}(T - S)\bigr)$，其中
$\vertiii{T^{-1}(T-S)} \leq \vertiii{T^{-1}}\,\vertiii{T - S} =
\theta < 1$：由\cref{prop:b3:banach:neumann}，$S$ 可逆，
且 $S^{-1} = \sum_{n\geq0}\bigl(T^{-1}(T -
S)\bigr)^nT^{-1}$，从而
\[
\vertiii{S^{-1} - T^{-1}} \leq
\sum_{n\geq1}\theta^n\,\vertiii{T^{-1}}
= \frac{\vertiii{T^{-1}}^2\,\vertiii{T - S}}{1 - \theta}.
\]
对线性方程组：$x_T = T^{-1}b$ 与 $x_S = S^{-1}b$
之差至多为上述界乘以 $\norm b$ --- 可逆系统的小
扰动仍唯一可解，
解对算子呈利普希茨依赖。
\end{solution}

\begin{solution}{exo:b3:banach:3}
(a) $\ell^1$：设 $(x^{(k)})$ 为柯西列。每个分量都是
柯西列（$\abs{x^{(k)}_n - x^{(l)}_n} \leq \norm{x^{(k)} -
x^{(l)}}_1$），记 $x_n = \lim_kx^{(k)}_n$。给定 $\varepsilon$，
当 $k, l \geq K$ 时，$\sum_{n \leq N}\abs{x^{(k)}_n - x^{(l)}_n}
\leq \varepsilon$ 对每个 $N$ 成立；令 $l \to \infty$，再令 $N \to
\infty$，得 $\norm{x^{(k)} - x}_1 \leq \varepsilon$，且 $x =
x^{(k)} - (x^{(k)} - x) \in \ell^1$。对 $\ell^\infty$，
$\norm\cdot_\infty$ 意义下的柯西列即一致柯西列，一致收敛于
一个有界数列。$c_0$ 在 $\ell^\infty$ 中是闭的：若
$x^{(k)} \to x$ 一致收敛，且 $x^{(k)}_n \to_n 0$，则
$\abs{x_n} \leq \norm{x - x^{(k)}}_\infty + \abs{x^{(k)}_n}$
得 $\limsup_n\abs{x_n} \leq \varepsilon$，故 $x \in c_0$；
巴拿赫空间的闭子空间仍是巴拿赫空间。有限序列：
它们的闭包包含每个 $x \in c_0$（截断收敛：
$\sup_{n>N}\abs{x_n} \to 0$），且包含于闭集
$c_0$ 中。

(b) 由齐次性，不妨设 $\norm x_p = 1$，则 $\abs{x_n} \leq
1$（对一切 $n$），从而 $\abs{x_n}^q \leq \abs{x_n}^p$，且 $\norm
x_q \leq 1 = \norm x_p$；当 $q = \infty$ 时，$\abs{x_n} \leq
\norm x_p$ 直接成立。严格性：取 $x_n = n^{-\alpha}$，其中
$\frac1q < \alpha \leq \frac1p$，则它属于 $\ell^q \setminus
\ell^p$（黎曼级数）。
\end{solution}

\begin{solution}{exo:b3:banach:4}
（$\leq$）若 $\norm f \leq 1$ 且 $f\restriction_F = 0$，则对每个 $y \in F$，
$\abs{f(x)} = \abs{f(x - y)} \leq \norm{x -
y}$；取下确界。（$\geq$，可取到）
由\cref{cor:b3:banach:hbnormed}(3) 得 $f$，使
$f\restriction_F = 0$，$\norm f \leq 1$，$f(x) = d(x, F)$，
此时上确界即为最大值。由此，$F \subseteq \bigcap\{\ker
f : f\restriction_F = 0\}$ 显然成立，而点 $x \notin F$
被上述泛函排除在交集之外
（$f(x) = d(x,F) > 0$，因 $F$ 为闭集）。
\end{solution}

\begin{solution}{exo:b3:banach:5}
对 $y \in \ell^\infty$，$\abs{\Lambda_y(x)} =
\abs{\sum x_ny_n} \leq \norm y_\infty\norm x_1$，故
$\norm{\Lambda_y} \leq \norm y_\infty$；取 $e_n$ 检验：
$\abs{y_n} = \abs{\Lambda_y(e_n)} \leq \norm{\Lambda_y}$：
等号成立。反之，给定 $\Lambda \in (\ell^1)'$，令 $y_n =
\Lambda(e_n)$：$\abs{y_n} \leq \norm\Lambda$，故 $y \in
\ell^\infty$；$\Lambda$ 与 $\Lambda_y$ 在有限
序列上一致，而有限序列在 $\ell^1$ 中稠密（$\norm{x - \sum_{n\leq
N}x_ne_n}_1 = \sum_{n>N}\abs{x_n} \to 0$），故 $\Lambda =
\Lambda_y$。映射 $y \mapsto \Lambda_y$ 是线性、保距的，
且为满射。对 $(\ell^\infty)'$，同样的开头给出一个数列
$y_n = \Lambda(e_n)$，但有限序列\emph{不}
在 $\ell^\infty$ 中稠密
（常数列 $\mathbf 1$ 与所有有限序列的距离均为 $1$
），故 $\Lambda$ 不由 $y_n$ 决定 ---
事实上 $(\ell^\infty)' \neq \ell^1$
（见\cref{pb:b3:banach:1}）。
\end{solution}

\begin{solution}{exo:b3:banach:6}
对每个固定的 $x$，$y \mapsto B(x, y)$ 是连续线性的：
$\sup_{\norm y \leq 1}\norm{B(x,y)} < \infty$。因此该族
$\{B(\cdot, y) : \norm y \leq 1\} \subseteq \mathcal L(E, G)$
（每个成员都连续，由关于 $x$ 的连续性得）逐点
有界于巴拿赫空间 $E$ 上：由巴拿赫–斯坦豪斯
（见\cref{thm:b3:banach:banachsteinhaus}）得存在 $C$，使得
$\norm{B(x,y)} \leq C\norm x$ 对一切 $\norm y \leq 1$ 成立；
再由 $y$ 的齐次性即得 $\norm{B(x,y)} \leq
C\norm x\norm y$。
\end{solution}

\begin{solution}{exo:b3:banach:7}
(a) $\norm f_1 \leq \norm f_\infty$：恒等映射有界，
且为双射。其逆映射无界：$f_n(x) = x^n$ 满足
$\norm{f_n}_1 = \frac1{n+1} \to 0$，但 $\norm{f_n}_\infty = 1$。
这与\cref{cor:b3:banach:equivnorms}并无矛盾：$(\mathcal
C(\intcc01), \norm\cdot_1)$ 是\emph{不}完备的
（见\cref{exo:b3:complete:1}）；该推论要求
两端均完备。

(b) 在有限数列构成的空间 $E_0$ 上（它在 $\ell^2$ 中稠密），
$\varphi(x) = \sum_n n\,x_n$ 线性且无界
（$\varphi(e_n) = n$，而 $\norm{e_n}_2 = 1$）。闭图像
定理不适用：$E_0$ 不完备 --- 且 $\varphi$
没有到 $\ell^2$ 上的连续延拓，这说明缺少一致连续性时，
仅靠稠密性是不够的
（见\cref{thm:b3:complete:extension}）。
\end{solution}

\begin{solution}{exo:b3:banach:8}
验证闭图像假设。设 $x_k \to x$ 且
$Tx_k \to z$（在 $\ell^2$ 中）。对每个 $y$：
\[
\langle z, y\rangle = \lim_k\langle Tx_k, y\rangle
= \lim_k \langle x_k, Ty\rangle = \langle x, Ty\rangle
= \langle Tx, y\rangle,
\]
利用内积在每个变元上的连续性，
（柯西–施瓦茨）及两次对称性。于是 $z - Tx$ 与
每个 $y$ 正交，特别地与自身正交：$z = Tx$。图像
是闭的，而 $\ell^2$ 是巴拿赫空间，故 $T$ 有界
（见\cref{thm:b3:banach:closedgraph}）。因此，对称算子
若定义在 \emph{整个} $\ell^2$ 上，则自动有界；
真正无界的对称算子（位置算子、动量算子、
哈密顿量）必须定义在真稠密子空间上。
\end{solution}

\begin{solution}{exo:b3:banach:9}
首先，$\norm{\Lambda_n} = \sum_i\abs{w_{i,n}}$：$\leq$ 为
三角不等式；$\geq$ 通过检验分段线性函数 $f$
满足 $\norm f_\infty \leq 1$ 且 $f(x_{i,n}) =
\operatorname{sign}(w_{i,n})$（在
有限个结点之间作线性插值；结点重合处符号一致）。

（$\Rightarrow$）若对每个 $f$ 逐点收敛，则得 (i)，
且逐点有界，由巴拿赫–施泰因豪斯
（见\cref{thm:b3:banach:banachsteinhaus}）在巴拿赫空间 $\mathcal
C(\intcc01)$ 上得 (ii)。

（$\Leftarrow$）设 $M = \sup_n\norm{\Lambda_n} + 1$。给定 $f$
及 $\varepsilon$，取多项式 $P$ 使 $\norm{f -
P}_\infty < \varepsilon/(2M)$
（见\cref{cor:b3:complete:weierstrass}）；于是
\[
\Bigl|\Lambda_n f - \int_0^1 f\Bigr|
\leq \abs{\Lambda_n(f - P)} + \Bigl|\Lambda_nP -
\int_0^1P\Bigr| + \Bigl|\int_0^1(P - f)\Bigr|
\leq \varepsilon + \Bigl|\Lambda_nP - \int_0^1P\Bigr| \to
\varepsilon .
\]
权重非负，且对常数精确：$\sum_i\abs{w_{i,n}}
= \sum_iw_{i,n} = \Lambda_n(\mathbf 1) = \int_0^1 1 = 1$，对于
对常数精确的求积公式 --- (ii) 以常数 $1$ 成立。
\end{solution}

\begin{solution}{exo:b3:banach:10}
$J$ 的线性性显然；$\norm{J(x)} = \sup_{\norm f \leq
1}\abs{f(x)} = \norm x$，由
\cref{cor:b3:banach:hbnormed}(2)。若 $\dim E = n$，则 $\dim E' =
n$（基给出坐标泛函），故 $\dim E'' = n$，
而单射（保距）的 $J$ 为满射。可分性：设
$(f_n)$ 在 $E'$ 中稠密，并取 $\norm{x_n} = 1$ 使
$\abs{f_n(x_n)} \geq \frac12\norm{f_n}$。令 $F =
\overline{\operatorname{Vect}}(x_n)$；若 $F \neq E$，取 $g
\in E'$，$g \neq 0$，且在 $F$ 上为零
（见\cref{cor:b3:banach:hbnormed}(3)）；取 $f_{n_k} \to g$：
\[
\norm{f_{n_k} - g} \geq \abs{(f_{n_k} - g)(x_{n_k})}
= \abs{f_{n_k}(x_{n_k})} \geq \tfrac12\norm{f_{n_k}}
\geq \tfrac12\bigl(\norm g - \norm{f_{n_k} - g}\bigr),
\]
故 $\norm{f_{n_k} - g} \geq \frac13\norm g > 0$，矛盾。
因此 $F = E$，以有理数（或 $\Q + \iu\Q$）为系数的
$x_n$ 的线性组合构成一个可数稠密集。
\end{solution}

\begin{solution}{exo:b3:banach:11}
(a) 良定义：$d(x, F)$ 只依赖于 $\bar x$
（将 $x$ 加上 $F$ 中元素不改变距离）。
$\norm\cdot$ 的齐次性与三角不等式
经下确界传递。正定性需要闭性：$\norm{\bar
x} = 0$ 意味着 $d(x, F) = 0$，即 $x \in \bar F = F$，即
$\bar x = 0$。$\vertiii\pi \leq 1$：$\norm{\bar x} \leq
\norm x$。开球映为开球：若 $\norm{\bar x} < 1$，
则某个代表元满足 $\norm{x - y} < 1$；反之
$\pi(B_E(0,1)) \subseteq B_{E/F}(0,1)$，由范数不等式
可得 --- 故 $\pi$ 为开映射，即开映射定理的
典型情形。

(b) 设 $\sum_k\norm{\bar x_k} < \infty$，并逐项取提升，
使 $\norm{x_k} \leq \norm{\bar x_k} + 2^{-k}$，则
$\sum\norm{x_k} < \infty$，故 $s = \sum_kx_k$ 在
巴拿赫空间 $E$ 中收敛（见\cref{exo:b3:complete:1}(b)），由
$\pi$ 的连续性得 $\sum_k\bar x_k = \bar s$：商空间 $E/F$ 中每个绝对
收敛级数都收敛，这与完备性等价
（同一练习）。

(c) 映射 $\lambda(x) = \lim_nx_n$ 是线性映射 $c \to K$，
它恰在 $c_0$ 上为零，因此诱导出线性双射
$c/c_0 \to K$。等距性：$d(x, c_0) = \abs{\lambda(x)}$ ---
$\leq$：从 $x$ 中减去数列 $x - \lambda(x)\mathbf
1 \in c_0$，余下 $\lambda(x)\mathbf 1$，其范数为
$\abs{\lambda(x)}$；$\geq$：对 $y \in c_0$，$\norm{x -
y}_\infty \geq \limsup_n\abs{x_n - y_n} = \abs{\lambda(x)}$。
\end{solution}

\begin{solution}{exo:b3:banach:12}
(a) $W = \ker P$ 为闭集（$0$ 在一个
连续映射下的原像）；$V = \operatorname{im}P = \ker(I - P)$
（事实上，$Px = x$ 当且仅当 $x \in \operatorname{im}P$，利用 $P^2 =
P$），同样为闭集。每个 $x$ 都可写成 $Px + (x - Px)$，
其中 $Px \in V$，$x - Px \in W$，且 $V \cap W = 0$（$x = Px
= 0$）：$E = V \oplus W$。

(b) 图像论证：设 $x_n \to x$ 且 $Px_n \to z$。
于是 $z \in V$（$V$ 闭，$Px_n \in V$），且 $x_n - Px_n \to
x - z \in W$（$W$ 闭）。因此 $x = z + (x - z)$，其中 $z \in
V$，$x - z \in W$；由分解的唯一性，$z =
Px$。$P$ 的图像闭，$E$ 是巴拿赫空间，因此 $P$
有界（\cref{thm:b3:banach:closedgraph}）。

(c) (a) 与 (b) 合起来即为等价性。在希尔伯特
空间中，每个闭的 $V$ 都有闭补空间 $V^\perp$
（\cref{ch:b3:hilbert}）：每个闭子空间都可补。
一般巴拿赫空间中此结论不成立 --- $c_0$
在 $\ell^\infty$ 中没有闭补（菲利普斯
定理，超出本课程工具范围），因此有界投影是一种特权，
而闭图像定理恰好就是
把几何分解转化为有界算子的
一套记账工具。
\end{solution}

\begin{solution}{pb:b3:banach:1}
\textbf{1.} 对 $a, b > 0$，由 $\log$ 的凹性，
$\log\bigl(\tfrac{a^p}p + \tfrac{b^q}q\bigr) \geq \tfrac1p\log
a^p + \tfrac1q\log b^q = \log(ab)$，取指数即得。（若 $ab = 0$
则不等式显然成立。）等号当且仅当 $a^p = b^q$。

\textbf{2.} 不妨设 $\norm x_p = \norm y_q = 1$
（齐次性；零的情形平凡）。于是
\[
\abs{\langle x, y\rangle} \leq \sum_n\abs{x_n}\abs{y_n}
\leq \sum_n\Bigl(\frac{\abs{x_n}^p}p +
\frac{\abs{y_n}^q}q\Bigr) = \frac1p + \frac1q = 1 =
\norm x_p\norm y_q .
\]
等号成立要求 $\abs{x_n}^p = \abs{y_n}^q$ 对一切 $n$
都成立（杨氏不等式的等号条件），并要求
$x_ny_n$ 的相位对齐。

\textbf{3.} $\abs{x_n + y_n}^p \leq \bigl(\abs{x_n} +
\abs{y_n}\bigr)\abs{x_n + y_n}^{p-1}$；求和后应用
赫尔德不等式（$p$ 与 $q$ 配对，注意 $(p - 1)q = p$）：
\[
\norm{x + y}_p^p \leq \bigl(\norm x_p + \norm
y_p\bigr)\,\Bigl(\sum_n\abs{x_n + y_n}^{p}\Bigr)^{1/q}
= \bigl(\norm x_p + \norm y_p\bigr)\,\norm{x+y}_p^{p/q};
\]
若 $\norm{x + y}_p \neq 0$（否则平凡），两边除以
$\norm{x+y}_p^{p/q}$，并利用 $p - \frac pq = 1$。（先说明
$\norm{x+y}_p$ 有限：$\abs{x_n+y_n}^p \leq
2^p(\abs{x_n}^p + \abs{y_n}^p)$。）

\textbf{4.} 完备性：与 $\ell^1$ 的情形相同
（见\cref{exo:b3:banach:3}），逐坐标取极限，再加上
一致的尾部估计 $\sum_{n\leq N}\abs{x^{(k)}_n - x^{(l)}_n}^p
\leq \varepsilon^p$，先令 $l$、再令 $N$ 趋于无穷。
有限序列的稠密性：$\norm{x - \sum_{n\leq
N}x_ne_n}_p^p = \sum_{n > N}\abs{x_n}^p \to 0$。

\textbf{5.} 由赫尔德不等式，$\abs{\Lambda_y(x)} \leq \norm
x_p\norm y_q$，故 $\norm{\Lambda_y} \leq \norm y_q$。取检验序列，令
$x^{(N)}_n = \abs{y_n}^{q-1}\overline{\operatorname{sign}}(y_n)$
当 $n \leq N$ 时，之后为 $0$（其中 $\operatorname{sign}$ 为
单位模相位，使得 $x_ny_n = \abs{y_n}^q$）。于是
$\Lambda_y(x^{(N)}) = \sum_{n\leq N}\abs{y_n}^q$，且
$\norm{x^{(N)}}_p = \bigl(\sum_{n\leq
N}\abs{y_n}^{q}\bigr)^{1/p}$（因 $(q-1)p = q$），故
\[
\norm{\Lambda_y} \geq \Bigl(\sum_{n\leq
N}\abs{y_n}^q\Bigr)^{1 - 1/p} \xrightarrow[N\to\infty]{}
\norm y_q .
\]

\textbf{6.} 令 $y_n = \Lambda(e_n)$。沿用相同的检验
向量，$\sum_{n \leq N}\abs{y_n}^q = \Lambda(x^{(N)}) \leq
\norm\Lambda\,\bigl(\sum_{n\leq N}\abs{y_n}^q\bigr)^{1/p}$,
于是 $\bigl(\sum_{n\leq N}\abs{y_n}^q\bigr)^{1/q} \leq
\norm\Lambda$ 对每个 $N$ 成立，从而 $y \in \ell^q$，$\norm y_q \leq
\norm\Lambda$。泛函 $\Lambda$ 与 $\Lambda_y$ 在
稠密的有限序列（第 4 题）上取值相同：$\Lambda =
\Lambda_y$。结合第 5 题，$y \mapsto \Lambda_y$ 是一个
保距同构 $\ell^q \cong (\ell^p)'$。

\textbf{7.} 设 $\xi \in (\ell^p)''$。与
第 6 题的保距同构 $\ell^q \cong (\ell^p)'$ 复合，$\xi$ 便确定了
$(\ell^q)'$ 中的一个元素，由第 6 题（$p, q$ 互换）可知
它为 $\Lambda_z$，其中 $z \in \ell^p$ 唯一确定：对每个
$y \in \ell^q$，$\xi(\Lambda_y) = \sum_nz_ny_n$。另一方面，
$J(z)(\Lambda_y) = \Lambda_y(z) = \sum_ny_nz_n$，两者取值相同。
由于 $(\ell^p)'$ 的每个元素都是某个 $\Lambda_y$，
故 $\xi = J(z)$：$J$ 为满射 --- $\ell^p$ 自反。

\textbf{8.} 元素为有理数（实部与虚部）的有限序列
可数，且在 $\ell^p$（$p < \infty$）中稠密，并
在 $c_0$ 中稠密。在 $\ell^\infty$ 中，族 $\{\mathbf 1_A : A
\subseteq \N\}$ 不可数，且 $\norm{\mathbf 1_A -
\mathbf 1_B}_\infty = 1$（当 $A \neq B$ 时），球 $B(\mathbf
1_A, \frac12)$ 两两不交，而稠密集须与每个球都相交：
不存在可数稠密集。

\textbf{9.} 若 $\ell^1$ 自反，则 $(\ell^\infty)'
\cong ((\ell^1)')' = J(\ell^1)$ 将可分（作为可分的 $\ell^1$ 的
等距像）。由
\cref{exo:b3:banach:10}可知，对偶空间
$(\ell^\infty)'$ 的可分性将迫使 $\ell^\infty$ 可分 ---
与第 8 题矛盾。故 $\ell^1$ 不自反（且
$(\ell^\infty)'$ 严格大于 $\ell^1$，正如
第 11 题具体说明的那样）。

\textbf{10.} $p$ 的齐次性显然；次可加性：
平均是线性的，且 $\limsup(u_n + v_n) \leq \limsup u_n +
\limsup v_n$。在 $c$ 上：收敛数列的切萨罗平均
收敛于其极限，故 $p(x) = \lim x =
\mathrm{LIM}(x)$ 于此成立；特别地，$\mathrm{LIM} \leq p$ 在
$c$ 上成立。由哈恩–巴拿赫（\cref{thm:b3:banach:hahnbanach}）
$\mathrm{LIM}$ 延拓到 $\ell^\infty_\R$ 上，并使 $\mathrm{LIM} \leq p$
全局成立。正性：若 $x \geq 0$，则 $-\mathrm{LIM}(x) =
\mathrm{LIM}(-x) \leq p(-x) = \limsup\text{平均}(-x) \leq 0$。
移位不变性：$x - Sx$ 的平均裂项相消为
$\frac{x_1 - x_{n+1}}n \to 0$，故 $p(\pm(x - Sx)) = 0$，且
$\mathrm{LIM}(x - Sx) = 0$。界：$\mathrm{LIM}(x) \leq p(x)
\leq \limsup x$（平均落后于上确界），并将此式
用于 $-x$，即得下界。

\textbf{11.} $e_n \in c_0$，故 $\mathrm{LIM}(e_n) = 0$
对所有 $n$。若 $\mathrm{LIM} = \Lambda_y$，$y \in \ell^1$，则
$y_n = \Lambda_y(e_n) = 0$，对所有 $n$，故 $\Lambda_y = 0$；但
$\mathrm{LIM}(\mathbf 1) = 1$。因此 $\mathrm{LIM} \in
(\ell^\infty)' \setminus \{\Lambda_y : y \in \ell^1\}$：
$(\ell^\infty)' \neq \ell^1$，再次成立。

\textbf{12.} 对 $x = (0,1,0,1,\dots)$：$x + Sx = \mathbf 1$，
故 $2\,\mathrm{LIM}(x) = \mathrm{LIM}(x) + \mathrm{LIM}(Sx) =
1$：$\mathrm{LIM}(x) = \frac12$。假设 $\varphi$ 是
极限的一个移位不变的乘性延拓：$x\cdot
Sx = 0$ 给出 $\varphi(x)\varphi(Sx) = \varphi(x)^2 = 0$，故
$\varphi(x) = 0$；但 $\varphi(x) + \varphi(Sx) =
\varphi(\mathbf 1) = 1$ 给出 $2\varphi(x) = 1$：矛盾。
平均与乘法不可能同时成立。

\textbf{13.} 设 $\norm{x^{(k)}}_p \leq M$。坐标
以 $M$ 为界：由对角线法得一子列
（仍记为 $x^{(k)}$），满足 $x^{(k)}_n \to x_n$，对每个
$n$。于是 $x \in \ell^p$：$\sum_{n\leq N}\abs{x_n}^p =
\lim_k\sum_{n\leq N}\abs{x^{(k)}_n}^p \leq M^p$，对一切 $N$。
弱收敛：对 $y \in \ell^q$ 及任意 $N$，
\[
\abs{\Lambda_y(x^{(k)} - x)} \leq
\Bigl|\sum_{n \leq N}(x^{(k)}_n - x_n)y_n\Bigr|
+ 2M\Bigl(\sum_{n>N}\abs{y_n}^q\Bigr)^{1/q},
\]
其中第一项趋于 $0$（当 $k \to \infty$ 时，因只涉及
有限多个坐标），第二项当 $N$ 充分大时很小
（对尾部用赫尔德不等式）：$\Lambda_y(x^{(k)}) \to \Lambda_y(x)$
对每个 $y$ 成立 --- 即弱收敛，因为每个泛函都是
某个 $\Lambda_y$（第 6 题）。在 $\ell^1$ 中此结论不成立：考虑
$(e_n)$，有界。任一子列 $(e_{n_k})$ 按坐标
收敛于 $0$，故其唯一可能的弱极限为 $0$；
但取 $y \in \ell^\infty$，使 $y_{n_k} =
(-1)^k$（其余为 $0$），则 $\Lambda_y(e_{n_k}) = (-1)^k$
发散。故无弱收敛子列：球的弱序列紧性
刻画了自反空间的情形。

\textbf{14.} 对 $y \in \ell^1$，$\Lambda_y(x) = \sum x_ny_n$
在 $c_0$ 上有定义，且 $\abs{\Lambda_y(x)} \leq \norm
x_\infty\norm y_1$，并取 $x^{(N)} =
(\operatorname{sign}y_1, \dots, \operatorname{sign}y_N, 0,
\dots) \in c_0$，则 $\Lambda_y(x^{(N)}) =
\sum_{n\leq N}\abs{y_n} \to \norm y_1$:
$\norm{\Lambda_y} = \norm y_1$。反之，对 $\Lambda \in
(c_0)'$，令 $y_n = \Lambda(e_n)$；同样的检验给出
$\sum_{n \leq N}\abs{y_n} = \Lambda(x^{(N)}) \leq
\norm\Lambda$，故 $y \in \ell^1$，且 $\Lambda = \Lambda_y$ 在
稠密的有限序列上成立，从而处处成立。
对偶链为 $(c_0)' = \ell^1$，$(\ell^1)' = \ell^\infty$
（\cref{exo:b3:banach:10}），$(\ell^\infty)' \supsetneq \ell^1$
（第 9--11 题）：自反性在第一步就已失效
--- $c_0'' = \ell^\infty \neq c_0$ --- 且再也没有恢复。

\textbf{15.} $J x^{(k)} \in E'' = (E')'$ 是
\emph{巴拿赫}空间 $E'$ 上的一族有界泛函（对偶
空间完备，\cref{prop:b3:banach:llcomplete}）；对每个
$\Lambda \in E'$，序列 $Jx^{(k)}(\Lambda) =
\Lambda(x^{(k)})$ 收敛，从而有界。
对 $E'$ 应用巴拿赫–施泰因豪斯定理，得 $\sup_k\norm{Jx^{(k)}} <
\infty$，且 $J$ 是等距的
（\cref{rem:b3:banach:bidual}）：$\sup_k\norm{x^{(k)}} <
\infty$.

\textbf{16.} （$\Rightarrow$）有界性即第 15 题；
坐标即为泛函 $\Lambda_{e_n}$。
（$\Leftarrow$）设 $M = \sup_k\norm{x^{(k)}}_p$，$y \in
\ell^q$，$\varepsilon > 0$；取 $N$ 使
$\bigl(\sum_{n>N}\abs{y_n}^q\bigr)^{1/q} < \varepsilon$。于是
\[
\abs{\Lambda_y(x^{(k)} - x)} \leq
\sum_{n\leq N}\abs{x^{(k)}_n - x_n}\,\abs{y_n} +
\norm{x^{(k)} - x}_p\,\varepsilon
\leq \sum_{n\leq N}\abs{x^{(k)}_n - x_n}\abs{y_n} +
(M + \norm x_p)\,\varepsilon,
\]
而有限和趋于 $0$：$\limsup \leq (M + \norm
x_p)\varepsilon$ 对每个 $\varepsilon$ 成立。（$x \in
\ell^p$ 且 $\norm x_p \leq M$，由法图型
的有限截断估计得出：$\sum_{n \leq N}\abs{x_n}^p =
\lim_k\sum_{n\leq N}\abs{x^{(k)}_n}^p \leq M^p$。）对 $e_k$
在 $\ell^2$ 中：有界，按坐标 $\to 0$，故 $e_k
\rightharpoonup 0$，但 $\norm{e_k} = 1$：单位质量
逃逸至无穷大指标，对每个固定的
泛函都不可见。

\textbf{17.} 取 $\Lambda$，满足 $\norm\Lambda = 1$ 且
$\Lambda(x) = \norm x$（\cref{cor:b3:banach:hbnormed}）。则
$\norm x = \lim\Lambda(x^{(k)}) \leq
\liminf\norm\Lambda\,\norm{x^{(k)}} =
\liminf\norm{x^{(k)}}$。（取 $e_k \rightharpoonup 0$：$0
\leq \liminf 1$，且不等式可以是严格的。）

\textbf{18.} 在 $\ell^2$ 中，$\norm{x^{(k)} - x}_2^2 =
\norm{x^{(k)}}_2^2 - 2\operatorname{Re}\langle x^{(k)},
x\rangle + \norm x_2^2$（实情形：$-2\langle x^{(k)},
x\rangle$）。弱收敛应用于泛函
$\Lambda_x$ 得 $\langle x^{(k)}, x\rangle \to \norm
x_2^2$，又由假设范数收敛，故右端
趋于 $\norm x^2 - 2\norm x^2 + \norm x^2 = 0$。
反例（不要求范数收敛）：$e_k
\rightharpoonup 0$，$\norm{e_k - 0} = 1 \not\to 0$。

\textbf{19.} 坐标收敛：取泛函
$\Lambda_{e_n} \in (\ell^1)' = \ell^\infty$ ($y = e_n$).
构造：设已选定 $k_{j-1}, N_{j-1}$，取 $k_j >
k_{j-1}$ 足够大，使得 $\sum_{n \leq
N_{j-1}}\abs{x^{(k_j)}_n} < \frac\delta{20}$（有限多个
坐标，每个都 $\to 0$），再取 $N_j > N_{j-1}$ 足够大，
使尾和 $\sum_{n > N_j}\abs{x^{(k_j)}_n} <
\frac\delta{20}$（即定义
$\norm{x^{(k_j)}}_1$ 的级数收敛）。区块 $B_j = \intoc{N_{j-1}}{N_j}$
则该区块承载除 $\frac\delta{10}$ 外的全部
$x^{(k_j)}$ 的质量。

\textbf{20.} 取如上定义的 $y$（$\abs{y_n} \leq 1$
处处成立）：
\[
\langle x^{(k_j)}, y\rangle
= \sum_{n \in B_j}\abs{x^{(k_j)}_n}
+ \sum_{n \notin B_j}x^{(k_j)}_ny_n
\geq \Bigl(\norm{x^{(k_j)}}_1 - \frac\delta{10}\Bigr) -
\frac\delta{10}
\geq \delta - \frac{2\delta}{10} = \frac{4\delta}5 > 0,
\]
中间的不等式成立，因为区块之外的质量至
多为 $\frac\delta{10}$（题~19）。但 $y \in
\ell^\infty = (\ell^1)'$ 且 $x^{(k)} \rightharpoonup 0$
迫使 $\langle x^{(k_j)}, y\rangle \to 0$，矛盾。
因此 $\ell^1$ 中的弱零序列依范数趋于零，由
平移可知弱收敛序列依范数收敛，这就是舒尔
定理。

\textbf{21.} 若 $(e_k)$ 的某子列弱收敛，其极限
须为 $0$（比较坐标），由舒尔定理得 $\norm{e_{k_j}}_1
\to 0$ --- 但范数都是 $1$。故
不存在弱收敛子列，与题~13 中手算所得一致。这并非
悖论：舒尔定理说明序列无法区分弱
拓扑与 $\ell^1$ 中的范数拓扑（拓扑本身
确有不同 --- 弱邻域从来都不是范数有界的），而
单位球的弱序列紧性则是另一种
更强的性质，与自反性等价
（埃伯莱因–什穆利安定理超出本书工具；其中失效的一侧，
至少我们已经证明）。

\textbf{22.} \emph{总表。}
\begin{center}
\begin{tabular}{l|l|c|c|c}
$E$ & $E'$ & 可分 & 自反 & 弱序列紧单位球\\
\hline
$c_0$ & $\ell^1$ & 是 & 否 & 否（$e_1{+}\dots{+}e_k$）\\
$\ell^1$ & $\ell^\infty$ & 是 & 否 & 否（$e_k$，题~21）\\
$\ell^p$ & $\ell^q$ & 是 & 是 & 是（题~13）\\
$\ell^\infty$ & $\supsetneq\ell^1$ & 否 & 否 & 否\\
\end{tabular}
\end{center}
特征：$c_0$ --- 其二重对偶为 $\ell^\infty$：首个
非自反的情形（题~14）；$\ell^1$ --- 舒尔
性质（题~20）；$\ell^p$ --- 自反性与弱
紧性（题~7、13）；$\ell^\infty$ ---
不可分性与巴拿赫极限：没有任何序列
能表示的泛函（题~8、10--11）。同一族空间，
四个不同的世界。

\textbf{23.} 设 $(f_k) \subseteq F$ 且 $\norm{x - f_k}
\to d$。则 $\norm{f_k} \leq \norm x + \sup_k\norm{x -
f_k}$：有界，由题~13 知存在子列 $f_{k_j}
\rightharpoonup f$。若 $f \notin F$，则 $\delta =
\operatorname{dist}(f, F) > 0$（$F$ 闭）；在 $F \oplus \R
f$ 上的线性型 $\lambda(g + tf) = t$ 满足
$\abs{\lambda(u)} \leq \norm u/\delta$（因为 $\norm{g +
tf} \geq \abs t\,\delta$），由哈恩--巴拿赫定理把它延拓为
$\Lambda \in (\ell^p)'$，且 $\Lambda\restriction_F = 0$，
$\Lambda(f) = 1$；但于是 $0 = \Lambda(f_{k_j}) \to
\Lambda(f) = 1$：矛盾。故 $f \in F$，且 $x -
f_{k_j} \rightharpoonup x - f$ 给出，由题~17 可得
\[
d \leq \norm{x - f} \leq \liminf_j\,\norm{x - f_{k_j}} = d :
\]
距离在 $f$ 处取到。在 $c_0$ 中：$\abs{\Lambda(x)}
\leq \sum_n2^{-n}\abs{x_n} < \norm x_\infty$，对每个 $x
\neq 0$（非零的趋零序列不可能满足 $\abs{x_n} =
\norm x_\infty$ 对所有 $n$ 都成立），而截断序列 $(1,
\dots, 1, 0, \dots)$ 给出 $\Lambda = 1 - 2^{-N} \to 1$：于是
$\norm\Lambda = 1$，从不取到。距离公式：对 $f
\in \ker\Lambda$，$\abs{\Lambda(x)} = \abs{\Lambda(x - f)}
\leq \norm{x - f}$，故 $\operatorname{dist}(x, \ker\Lambda)
\geq \abs{\Lambda(x)}$；反之，对单位球中的 $u$
使 $\Lambda(u) \geq 1 - \varepsilon$，取向量 $f = x -
\frac{\Lambda(x)}{\Lambda(u)}u$ 属于 $\ker\Lambda$，且
$\norm{x - f} \leq \frac{\abs{\Lambda(x)}}{1 - \varepsilon}$：
等号成立。若某个 $f \in \ker\Lambda$ 达到此值，则 $z = x -
f$ 满足 $\abs{\Lambda(z)} = \abs{\Lambda(x)} =
\norm z \neq 0$，于是 $\Lambda$ 会在
$z/\norm z$ 处取得其范数，不可能。$c_0$ 的一个闭超平面，其上
处处没有最近点 --- 自反性绝非装饰。

\textbf{24.} 令 $y_1 = x^{(1)}$。给定 $y_1, \dots, y_j$，
每个映射 $k \mapsto \langle y_i, x^{(k)}\rangle$ 都趋于 $0$，
（$y_i \in \ell^2 = (\ell^2)'$），故存在超过前一指标的 $k_{j+1}$，
使 $\abs{\langle y_i,
x^{(k_{j+1})}\rangle} \leq \frac1{j+1}$ 对 $i = 1, \dots,
j$ 成立；记此选取为 $y_{j+1}$。于是
\[
\Bigl\lVert\sum_{j=1}^my_j\Bigr\rVert_2^2
= \sum_{j=1}^m\norm{y_j}_2^2
+ 2\sum_{j=2}^m\sum_{i<j}\langle y_i, y_j\rangle
\leq mC^2 + 2\sum_{j=2}^m\frac{j-1}j
\leq mC^2 + 2m,
\]
两边除以 $m^2$，得 $\norm{\frac1m\sum_jy_j}_2^2 \leq
\frac{C^2 + 2}m \to 0$。（若弱极限 $x \neq 0$，
则对 $x^{(k)} - x$ 应用此结论。）在标准正交列 $(e_k)$ 上，连
子列也无需抽取：$\norm{\frac1m(e_1 + \dots +
e_m)}_2 = \frac{\sqrt m}m = \frac1{\sqrt m}$。平均
把弱收敛化为依范数收敛，此即
$\ell^2$ 的巴拿赫--萨克斯性质。

\textbf{25.} $A_mx = \frac1m\sum_{i=0}^{m-1}S^ix$，故
由线性与平移不变性得 $L(A_mx) = L(x)$。对任意
有界 $u$ 及 $\varepsilon > 0$，取 $N$，使 $u_n \leq
\limsup u + \varepsilon$（$n \geq N$）；由正性，
$(\limsup u + \varepsilon)\mathbf 1 - S^Nu \geq 0$ 及
$L(S^Nu) = L(u)$，可得 $L(u) \leq \limsup u + \varepsilon$，
同理可得 $L(u) \geq \liminf u - \varepsilon$，因此
$\liminf A_mx \leq L(x) \leq \limsup A_mx$ 对每个 $m$ 成立。若
$x$ 为 $T$ 周期的，则 $A_Tx$ 为常数列，其值等于
周期均值 $\mu$：$L(x) = \mu$，对\emph{每一个}巴拿赫
极限 --- $\frac13$ 对应 $(1,0,0,\dots)$，$\frac12$ 对应
$(0,1,0,1,\dots)$ 的情形，如第 12 题所示。对于块序列：
取 $N = 3^j$，$j$ 为偶数时，最后一块全为 1，故
切萨罗均值 $\geq \frac{3^j - 3^{j-1}}{3^j} = \frac23$；
取 $N = 3^j$，$j$ 为奇数时，所有 1 均位于 $\intoc0
{3^{j-1}}$ 中，故均值 $\leq \frac13$。因此 $p(x) \geq
\frac23$，且 $-p(-x) = \liminf_n\frac{x_1 + \dots + x_n}n
\leq \frac13$。在 $M = c \oplus \R x$ 上定义 $\Lambda_+(y +
tx) = \lim y + t\,p(x)$。关于 $p$ 的控制性：当 $t > 0$ 时，
由次可加性，$p(tx) \leq p(y + tx) + p(-y)$，即
$p(y + tx) \geq t\,p(x) + \lim y$（注意 $p(\pm y) = \pm\lim
y$，对 $y \in c$：收敛数列的切萨罗均值
当 $t = -s < 0$ 时，$p(y) \leq p(y -
sx) + p(sx)$，得 $p(y - sx) \geq \lim y - s\,p(x)$；当 $t
= 0$ 时等号成立。于是 $\Lambda_+ \leq p$ 在 $M$ 上，
由哈恩--巴拿赫定理，将其延拓为 $L_+ \leq p$，定义在 $\ell^\infty_\R$ 上，
它恰为第 10 题中的巴拿赫极限（
由 $p$ 的控制性可推出正性、平移不变性，以及
$\lim$ 在 $c$ 上的取值），$L_+(x) = p(x) \geq \frac23$。同样的
计算，取 $\Lambda_-(y + tx) = \lim y - t\,p(-x)$，
（利用 $p(-sx) \leq p(y - sx) + p(-y)$，对情形 $t = -s
< 0$）得巴拿赫极限 $L_-$，满足 $L_-(x) = -p(-x) \leq
\frac13$。两个巴拿赫极限，同一数列，两个取值。
在周期（以及更一般地，\emph{几乎
收敛}）数列的情形之外，巴拿赫极限是一种真正的选择。
\end{solution}
